% Identity-Aware Merge for Structured Data:
% Field-Wise Pushout with Separated Identity and Location
%
% Draft for arXiv cs.SE

\documentclass[11pt,a4paper]{article}
\usepackage[margin=1in]{geometry}
\usepackage{amsmath,amssymb,amsthm}
\usepackage{booktabs}
\usepackage{hyperref}
\usepackage{enumitem}

\newtheorem{theorem}{Theorem}
\newtheorem{lemma}[theorem]{Lemma}
\newtheorem{corollary}[theorem]{Corollary}
\newtheorem{definition}[theorem]{Definition}
\newtheorem{remark}[theorem]{Remark}

\newcommand{\Id}{\mathsf{Id}}
\newcommand{\Field}{\mathsf{Field}}
\newcommand{\Node}{\mathsf{Node}}
\newcommand{\Section}{\mathsf{Section}}
\newcommand{\Patch}{\mathsf{Patch}}

\title{Merge for Structured Data as a Sheaf on the Tree Space:\\
Referential Integrity via Sheafification}

\author{
Jihao Yang\thanks{AstraZeneca R\&D China, Shanghai, China}
}

\date{\today}

\begin{document}
\maketitle

\begin{abstract}
Three-way merge for tree-structured data (JSON, YAML, XML) is commonly framed
as a per-field pushout on a discrete product of stalks --- one factor per
(identity, field) pair. We argue this framing is too weak: by placing the
topology on the \emph{tree} (the Alexandrov topology of the ancestry poset)
rather than on a discrete set, merge becomes the pushout in the sheaf category
$Sh(\Id)$, and sheafification enforces a \emph{referential-integrity
constraint} (a present node must have a present parent). This gives a new class
of merge obstruction --- a \emph{structural conflict} (a present child whose
parent was deleted on the other branch) --- that the discrete per-field model
is provably blind to.

Concretely, we (i) separate identity from location so that Move is a field
change, enabling Move+Modify to commute; (ii) define merge as the two-stage
pushout $\Sigma(m_0)$ --- pointwise per-field pushout $m_0$ in the presheaf
category, followed by sheafification $\Sigma$ that drops dangling references to
a fixpoint; (iii) prove a Strict-Refinement theorem: the sheaf conflict set
strictly contains the discrete one, the gap being exactly the structural
conflicts, non-empty iff the tree has depth $\geq 1$ and the two branches'
edits cross an ancestry edge inconsistently. We implement the algorithm in
Rust (open source, 78 tests) and verify the sheaf-consistency invariant and
the groupoid laws by property-based testing (2000 cases each).
\end{abstract}


% ============================================================
\section{Introduction}
% ============================================================

Version control systems merge concurrent edits by comparing two branches
against a common base. Git performs this comparison at the granularity of
\emph{text lines}: two edits conflict if they fall in the same 3-line
context window. Pijul performs it at the granularity of \emph{patches}
on a flat sequence, with commutation rules for adjacent lines.

Both approaches are designed for \emph{source code}---flat sequences of
text lines. But an increasing fraction of version-controlled data is
\emph{structured}: JSON configurations, YAML manifests, XML documents,
database schemas. These have a tree shape with stable identities (object
keys, element IDs, primary keys). For such data, two distinct defects of
line-based merge arise. First, \emph{false conflicts}: two developers
editing different fields of the same JSON object trigger a conflict
because their edits are textually adjacent, even though they are logically
disjoint. Second, and less noticed, \emph{silent structural violations}:
when two branches cross an ancestry edge inconsistently --- for instance,
one branch deletes a node while another moves a child under it ---
line-based merge (and, as we show, the discrete per-field merge) will
happily keep the child with a reference to a now-absent parent --- a
dangling reference that is not a section of any well-formed tree.

\paragraph{Contribution.}
We make four contributions:

\begin{enumerate}[noitemsep]
\item \textbf{The tree-space sheaf} (\S\ref{sec:model}): we place the
topology on the tree (Alexandrov topology of the ancestry poset) rather
than on a discrete key set. Sections of the sheaf satisfy a
referential-integrity constraint: a present node has a present parent.

\item \textbf{Two-stage sheaf merge} (\S\ref{sec:merge}): merge is the
pushout $\Sigma(m_0)$ in $Sh(\Id)$ --- a pointwise per-field pushout $m_0$
in the presheaf category, followed by sheafification $\Sigma$ that drops
dangling references to a fixpoint. Combined with identity-location
separation, Move+Modify merges cleanly because they touch different fields.

\item \textbf{Theorems} (\S\ref{sec:theory}): Sheaf Pushout Correctness,
Sheaf Consistency (the output is always a valid section), and Strict
Refinement over the discrete per-field merge --- the gap being exactly the
new \emph{structural conflict} class, invisible to any discrete model.

\item \textbf{Implementation} (\S\ref{sec:impl}): an open-source Rust
library (\texttt{tate}) with 78 tests; the sheaf-consistency invariant
and the patch groupoid laws are verified by proptest (2000 cases each).
\end{enumerate}


% ============================================================
\section{The Tree-Space Sheaf}
\label{sec:model}
% ============================================================

The decisive modelling choice is \emph{where the topology lives}. Prior
structured-merge work (and the naive reading of "per-field pushout") treats
the (identity, field) pairs as a \emph{discrete} product of stalks. We place
the topology on the tree itself.

\begin{definition}[Ancestry poset]
The identities $\Id$ of a tree carry the ancestry order $I \sqsubseteq J$
iff $I = J$ or $I$ is on the parent chain from the root to $J$. The root
is the unique least element.
\end{definition}

\begin{definition}[Tree space]
The \emph{tree space} is $\Id$ equipped with the Alexandrov topology of
$\sqsubseteq$: $U \subseteq \Id$ is open iff it is ancestor-closed,
\[
J \in U \wedge I \sqsubseteq J \;\Longrightarrow\; I \in U.
\]
Opens are ancestor-closed subtrees. The smallest open containing $I$ is
$\downarrow I = \{\text{root},\dots,\text{parent}(I), I\}$. The topology is
discrete iff the tree is flat ($\sqsubseteq$ is equality).
\end{definition}

\begin{definition}[Stalk]
$F_I = \{\bot\} \cup \Node_I$, where $\bot$ = absent and $\Node_I$ is the
set of node records $(parent, kind, label, text, attrs, order)$.
\end{definition}

\begin{definition}[Section]
For an open $U$, a \emph{section} $s \in F(U)$ is an assignment
$s : U \to \bigsqcup_I F_I$ with $s(I) \in F_I$, satisfying the
\emph{consistency condition} for every present $I \in U$:
\begin{equation}
\tag{C}
s(I).parent = \bot \;\;\vee\;\; \bigl(s(I).parent = J \wedge J \in U \wedge s(J) \neq \bot\bigr).
\label{eq:consistency}
\end{equation}
A present node has a present parent. Because $U$ is ancestor-closed,
$J \in U$ holds automatically when $I$ has parent $J$, so (\ref{eq:consistency})
is checkable within $U$. With the obvious restriction maps $F$ is a sheaf:
sections agreeing on overlaps glue uniquely (the consistency condition is
local to the ancestor chain $\downarrow I$). A \emph{global section}
$S \in F(\Id)$ is the formal counterpart of a consistent tree.
\end{definition}

\begin{remark}[Identity-location separation]
The $parent$ field is part of the node \emph{record}, not part of the key.
Identity (the map key) and location (the $parent$ value) are separated.
A Move changes $parent$ at a stable identity; a Modify changes $text$ or
$attrs$ at a stable identity. This is what lets Move+Modify commute
(\S\ref{sec:algebra}); it is orthogonal to the topology but composes with
it cleanly.
\end{remark}

\paragraph{Why the topology is load-bearing.}
On a discrete space every presheaf is a sheaf and condition (\ref{eq:consistency})
is vacuous (no ancestry relation). On the tree space, (\ref{eq:consistency})
is a genuine constraint, and --- as we show in \S\ref{sec:merge}---its
enforcement during merge surfaces a class of obstructions that no discrete
model can detect.


% ============================================================
\section{Diff, Apply, Invert, Compose}
\label{sec:algebra}
% ============================================================

\begin{definition}[Patch]
A \emph{patch} records the old and new node for each identity that changed:
\[
\Patch = \Id \rightharpoonup (\Node_{\perp} \times \Node_{\perp})
\]
where $\Node_{\perp} = \Node \cup \{\perp\}$ and $\perp$ represents absence
(the identity did not exist, or ceases to exist).
\end{definition}

\begin{definition}[Diff]
$\mathrm{diff}(S_1, S_2) = \{ I \mapsto (S_1(I), S_2(I)) \mid S_1(I) \neq
S_2(I) \}$.
\end{definition}

\begin{definition}[Apply]
$\mathrm{apply}(P, S)$: for each $I \in \mathrm{dom}(P)$, verify $S(I) = P(I).old$,
then set $S(I) \leftarrow P(I).new$ (or remove $I$ if $new = \perp$).
Fails if any precondition is not met.
\end{definition}

\begin{definition}[Invert]
$\mathrm{invert}(P) = \{ I \mapsto (P(I).new, P(I).old) \mid I \in \mathrm{dom}(P) \}$.
\end{definition}

\begin{definition}[Compose]
$\mathrm{compose}(P, Q)$: for each identity $I$ in $\mathrm{dom}(P) \cup \mathrm{dom}(Q)$,
\[
old = \begin{cases} P(I).old & \text{if } I \in \mathrm{dom}(P) \\ Q(I).old & \text{otherwise} \end{cases}
\quad
new = \begin{cases} Q(I).new & \text{if } I \in \mathrm{dom}(Q) \\ P(I).new & \text{otherwise} \end{cases}
\]
If $old = new$, the edit cancels (is omitted).
\end{definition}

These operations satisfy the groupoid laws (verified by property-based
testing with 2000 random cases each, see \S\ref{sec:impl}):
\begin{align}
\mathrm{apply}(\mathrm{diff}(a,b),\, a) &= b \label{eq:roundtrip}\\
\mathrm{apply}(\mathrm{invert}(P),\, \mathrm{apply}(P, S)) &= S \label{eq:invert}\\
\mathrm{apply}(\mathrm{compose}(P, Q),\, S) &= \mathrm{apply}(Q,\, \mathrm{apply}(P, S)) \label{eq:compose}\\
\mathrm{compose}(P,\, \mathrm{invert}(P)) &= \emptyset \label{eq:cancel}\\
\mathrm{compose}(P,\, \emptyset) &= P = \mathrm{compose}(\emptyset,\, P) \label{eq:identity}\\
\mathrm{compose}(\mathrm{compose}(P, Q),\, R) &= \mathrm{compose}(P,\, \mathrm{compose}(Q, R)) \label{eq:assoc}
\end{align}


% ============================================================
\section{Two-Stage Sheaf Merge}
\label{sec:merge}
% ============================================================

Given base $B$, ours $O$, theirs $T \in F(\Id)$, merge is the pushout of
the span $O \leftarrow B \rightarrow T$ in the category $Sh(\Id)$ of sheaves
on the tree space. In $Sh(\Id)$, colimits are obtained by sheafifying the
presheaf colimit. This yields a two-stage algorithm.

\paragraph{Stage 1: pointwise pushout $m_0$ (presheaf).}
Forget the topology and work in $PSh(\Id) = \prod_I \mathbf{Set}^{F_I}$.
At each identity $I$, take the pushout in $\mathbf{Set}$ of
$O(I) \leftarrow B(I) \rightarrow T(I)$ field-by-field:
\[
m_0(I).f = \begin{cases}
T(I).f & \text{if } O(I).f = B(I).f \\
O(I).f & \text{if } T(I).f = B(I).f \\
O(I).f & \text{if } O(I).f = T(I).f \\
\text{conflict} & \text{otherwise}
\end{cases}
\]
At a field conflict the node is marked \texttt{Field}-conflicted and takes
$O(I)$ as best-effort. This $m_0$ is the classical discrete per-field
pushout. It need not satisfy the consistency condition (\ref{eq:consistency}).

\paragraph{Stage 2: sheafification $\Sigma$.}
Let $\Sigma : PSh(\Id) \to Sh(\Id)$ be sheafification. Concretely, $\Sigma$
iterates to a fixpoint: drop every present node whose parent is absent.
Each dropped identity is a \texttt{Dangling} conflict. The merged section is
\[
M = \Sigma(m_0).
\]

\begin{definition}[Conflict set]
$C(B,O,T) \;=\; C_{\text{field}}(m_0) \;\uplus\; C_{\text{struct}}(\Sigma)$,
where $C_{\text{field}}$ collects Stage-1 per-$(I,f)$ pushout failures and
$C_{\text{struct}}$ collects the support of $\Sigma$'s correction (nodes
dropped for dangling parents).
\end{definition}

\paragraph{Move + Modify.}
If ours moves $I$ (changes $parent$) and theirs modifies $I$ (changes
$text$/$attrs$), then at $(I,parent)$: $O \neq B$, $T = B$ $\Rightarrow$
take $O$; at $(I,text)$: $O = B$, $T \neq B$ $\Rightarrow$ take $T$. No
field conflict, and the moved node's new parent is present by Move's
precondition, so Stage 2 does not fire. Move + Modify merges cleanly.

\paragraph{The new case: move-child + delete-target.}
Suppose ours deletes a node $Q$ while theirs moves $C$ so that
$parent(C)$ becomes $Q$. Both inputs are sections (the moved $C$ sits under
a present $Q$ in theirs; ours simply omits $Q$). Stage~1 takes the deletion
at identity $Q$ ($O{\neq}B$, $T{=}B$) and the move at $(C,parent)$
($O{=}B$, $T{\neq}B$) --- each single-sided, hence zero field conflicts.
But the merged $m_0$ has $C$ present with $parent(C){=}Q$ absent, violating
(\ref{eq:consistency}); $\Sigma$ drops $C$ and records a \texttt{Dangling}
conflict. \emph{The discrete per-field merge reports this as clean and
returns a non-section}; the sheaf merge flags it. This is the strict
refinement of \S\ref{sec:theory}.


% ============================================================
\section{Theorems and Proofs}
\label{sec:theory}
% ============================================================

\subsection{Patch Commutation (Stage 1)}

\begin{theorem}[Field Independence]
\label{thm:field}
Let $P$ and $Q$ be patches. If for every identity $I$, the fields touched
by $P$ at $I$ and by $Q$ at $I$ are disjoint, then
$\mathrm{compose}(P, Q) = \mathrm{compose}(Q, P)$.
\end{theorem}

\begin{proof}
At each $(I, f)$: if $f$ is touched only by $P$, both orders yield
$(P.\text{old}, P.\text{new})$; symmetric for $Q$ alone; if neither, no
edit. All $(I,f)$ agree.
\end{proof}

\begin{corollary}[Move-Modify Commutation]
$\mathrm{Move}(I, A{\to}B)$ and $\mathrm{Modify}(I, v{\to}w)$ commute.
\end{corollary}

\begin{proof}
Move touches $(I,parent)$, Modify touches $(I,text)$ or $(I,attrs)$;
$\{parent\} \cap (\{text\} \cup \{attrs\}) = \emptyset$. These theorems
live in Stage 1 and are unaffected by $\Sigma$: a well-formed Move's new
parent $B$ is present, so Stage 2 never fires on a Move.
\end{proof}

\subsection{Sheaf Pushout and Consistency}

\begin{theorem}[Sheaf Pushout Correctness]
\label{thm:pushout}
$\mathrm{merge}(B,O,T) = \Sigma(m_0)$, the pushout of $O \leftarrow B
\rightarrow T$ in $Sh(\Id)$, with best-effort $O$ chosen at each
\texttt{Field}-conflicted stalk.
\end{theorem}

\begin{proof}
In the category $Sh(\Id)$ of sheaves on the tree space, colimits are obtained
by sheafifying the presheaf colimit (apply $\Sigma$). The presheaf pushout is
the pointwise one $m_0$ (a product of stalks has pointwise colimits; each
factor $\mathbf{Set}$'s pushout is the four-way rule). The algorithm builds
$m_0$ then applies $\Sigma$ to a fixpoint.
\end{proof}

\begin{theorem}[Sheaf Consistency Invariant]
\label{thm:consistency}
For all $B, O, T \in F(\Id)$, the output $M = \mathrm{merge}(B,O,T)$ is a
global section: every present node has a present parent (or is the root).
\end{theorem}

\begin{proof}
$\Sigma$ is the fixpoint of ``drop present nodes with absent parents.''
The fixpoint exists (finite $\Id$, each pass removes $\geq 1$ node) and by
construction satisfies (\ref{eq:consistency}): otherwise the pass would
not have terminated. Verified by proptest (2000 random scenarios).
\end{proof}

\subsection{Strict Refinement over the Discrete Merge}

\begin{theorem}[Strict Refinement]
\label{thm:strict}
Let $C_{\text{sheaf}}$ be the conflict set of $\mathrm{merge}(B,O,T)$ and
$C_{\text{disc}}$ that of the discrete Stage-1 merge alone.
\begin{enumerate}[noitemsep]
\item[(a)] $C_{\text{disc}} \subseteq C_{\text{sheaf}}$ (every discrete
  conflict is a sheaf \texttt{Field} conflict);
\item[(b)] the containment is strict: $\exists\, B,O,T$ with
  $C_{\text{disc}} = \emptyset$ but $C_{\text{sheaf}} \neq \emptyset$.
\end{enumerate}
\end{theorem}

\begin{proof}
(a) Stage 1 is identical in both; every discrete conflict appears in
$C_{\text{sheaf}}$ as $C_{\text{field}}$.

(b) Witness (four identities). Let $\Id = \{\text{root}, P, Q, C\}$ with
$parent(P){=}parent(Q){=}\text{root}$ and $parent(C){=}P$ in the base. All
three inputs are global sections:
\[
B = \{\text{root}, P, Q, C\},\quad
O = \{\text{root}, P, C\},\quad
T = \{\text{root}, P, Q, C\text{ with } parent(C){=}Q\}.
\]
Ours deletes $Q$ ($C$ retains $parent(C){=}P$, still present $\Rightarrow$ a
section); theirs moves $C$ under $Q$ ($Q$ present $\Rightarrow$ a section).

\emph{Discrete Stage~1.} At identity $Q$: $O(Q){=}\bot$,
$B(Q){=}T(Q){=}Q$, so $T{=}B$ $\Rightarrow$ take $O{=}\bot$ (deletion wins);
$Q$ is absent. At $(C,parent)$: $O{=}B{=}P$, $T{=}Q$, so $O{=}B$
$\Rightarrow$ take $T{=}Q$ (move wins). Every other field agrees. No field
is double-changed, so $C_{\text{disc}} = \emptyset$. Yet the result has $C$
present with $parent(C){=}Q$ absent: not a section.
\emph{Sheaf.} $\Sigma$ drops $C$, recording \texttt{Dangling}@$C$ with
$\mathit{missing\_parent}{=}Q$; $C_{\text{sheaf}} \neq \emptyset$.
\end{proof}

\begin{remark}
This is the smallest witness: with three or fewer identities, any section
that deletes a node must also delete or reparent its descendants, so no
dangling reference can survive into $m_0$. A fourth identity is needed so
that the deletion (of $Q$) and the reparenting (of $C$ onto $Q$) live on
different branches. The class of such witnesses is exactly the support of
$\Sigma$'s correction --- non-empty iff the tree has depth $\geq 1$ and the
two branches' edits cross an ancestry edge inconsistently. On a flat
(discrete) tree, $\Sigma = \mathrm{id}$ and Theorem~\ref{thm:strict}
collapses to equality: the discrete merge \emph{is} the sheaf merge in that
degenerate case.
\end{remark}

\subsection{Clean Merge}

\begin{theorem}[Clean Merge]
$\mathrm{merge}(B,O,T)$ is conflict-free iff (i) for every $(I,f)$ where
both branches differ from base they agree, and (ii) $m_0$ already satisfies
(\ref{eq:consistency}).
\end{theorem}

\begin{proof}
(⟸) (i) ⇒ no Stage-1 conflicts; (ii) ⇒ $\Sigma(m_0)=m_0$ ⇒ no Stage-2
conflicts. (⟹) contrapositive of each.
\end{proof}


% ============================================================
\section{Implementation}
\label{sec:impl}
% ============================================================

We implement the model in Rust as the open-source library \texttt{tate}
(v3.0.0, MIT license, available on crates.io).

\paragraph{Data structures.} \texttt{Section} is a
\texttt{BTreeMap<String, Node>}, where \texttt{Node} stores
\texttt{parent: Option<String>}, \texttt{kind}, \texttt{text},
\texttt{attrs: Vec<(String,String)>}, and \texttt{order: usize}.

\paragraph{The two stages.} Stage~1 is \texttt{merge\_sections}' per-field
loop applying the four-way rule at each identity. Stage~2 is the
\texttt{sheafify} fixpoint: a pass \texttt{dangling\_refs} collects present
nodes whose \texttt{parent} is absent from the map, records each as a
\texttt{SectionConflict} of kind \texttt{Dangling}, removes it, and repeats
until none remain. The output is a consistent global section by
Theorem~\ref{thm:consistency}.

\paragraph{Conversion.} \texttt{TreeNode::to\_section()} flattens a tree
into a section (assigning each node its identity, or its kind for keyless
nodes); \texttt{Section::to\_tree()} rebuilds by walking parent pointers
from the root.

\paragraph{Property-based testing.} The groupoid laws, the field-wise
pushout law, and the sheaf-consistency invariant
(\texttt{law\_merged\_section\_has\_no\_dangling\_parents}) are verified
by \texttt{proptest} with 2000 randomly generated trees per law. Trees are
generated with globally unique identities, up to depth 4 and 32 nodes.

\paragraph{Repository.} \url{https://github.com/j-yang/tate}


% ============================================================
\section{Evaluation}
\label{sec:eval}
% ============================================================

The central empirical claim of the paper is the Strict-Refinement theorem
(Theorem~\ref{thm:strict}): the sheaf merge surfaces a class of
obstructions invisible to the discrete merge. We support it two ways.

\paragraph{Witness test.} The unit test
\texttt{move\_child\_delete\_target\_is\_structural\_conflict}
instantiates the theorem's witness ($\Id = \{\text{root}, P, Q, C\}$; ours
deletes $Q$, theirs moves $C$ under $Q$) and asserts (i) the sheaf merge
reports a \texttt{Dangling} conflict at $C$ with \texttt{missing\_parent = Q},
and (ii) $C$ is dropped from the merged section. The same input under the
discrete Stage-1 merge produces no conflict and a non-section result.

\paragraph{Universal invariant.} The proptest
\texttt{law\_merged\_section\_has\_no\_dangling\_parents} generates 2000
random $(B, O, T)$ triples (independent random trees sharing only the
root identity) and asserts the merged section never contains a present
node with an absent parent. This exercises Theorem~\ref{thm:consistency}
across the whole random input space.

\paragraph{What we do not claim.} We do not report a false-conflict rate
of line-based merge versus ours on random trees: such a number is an
artifact of the chosen serialization and diff backend, and a naive
line-proximity simulation inflates it arbitrarily. A real-world
evaluation on Git repositories with JSON/YAML is left as future work.


% ============================================================
\section{Related Work}
% ============================================================

\paragraph{Structured and semi-structured merge.} Apel et al.'s
semistructured merge~\cite{semistructured} diff-merges ASTs and reports
reductions in false conflicts versus line-based merge for source code.
Tools like 3DM (XML), \texttt{jd} and \texttt{jsondiffpatch} (JSON) perform
3-way merge at the granularity of object keys. These systems separate
identity from location in practice but, to our knowledge, do not formalise
merge as a sheaf on the tree space nor prove a structural-conflict class
that the discrete per-field model misses (Theorem~\ref{thm:strict}).

\paragraph{CRDTs.} Conflict-free Replicated Data Types~\cite{crdt}
(Automerge, Yjs) provide strong eventual consistency for collaborative
editing via operation commutation. They use stable node identifiers
(actor + counter) so that move is a position change, not delete+insert ---
the same identity-location split we adopt. Our Field-Independence theorem
(Theorem~\ref{thm:field}) recovers a form of their commutativity from
field disjointness rather than operation semantics. CRDT tree types
typically do not model the ancestry topology as a sheaf, and the
dangling-reference problem after concurrent delete-parent / modify-child
is handled operationally (e.g.\ tombstones) rather than as a categorical
obstruction.

\paragraph{Pijul.} Pijul~\cite{pijul} implements a patch theory for text
based on commutation of line-level operations on a flat sequence. Move in
Pijul changes the line's position (its key), so Move + Modify does not in
general commute. Our Move-Modify commutation (Corollary after
Theorem~\ref{thm:field}) is enabled by the identity-location split.

\paragraph{Git.} Git's recursive 3-way merge~\cite{git} operates on text
diffs and has no notion of identity: moving a JSON object key is a
delete+insert at two line ranges, and a delete-parent / modify-child pair
produces a silently dangling result with no conflict marker.

\paragraph{Sheaf theory and finite spaces.} Our formalism uses the standard
sheafification functor on a topological space~\cite{godement} and the
correspondence between finite posets and finite Alexandrov
spaces~\cite{mccord, barmak}. The non-trivial content is that the ancestry
topology of a tree is non-discrete as soon as the tree has depth $\geq 1$, so
sheafification is non-identity and surfaces a structural-conflict class
invisible to any model that treats the object set as discrete. On a flat tree
the topology is discrete, every presheaf is a sheaf, and our merge
specialises to the classical per-field pushout.


% ============================================================
\section{Limitations and Future Work}
% ============================================================

\paragraph{Keyless data.} Nodes without explicit identity (e.g., JSON
array items) need auto-assigned identities; same-kind keyless siblings
collide at one map key. Alignment heuristics (LCS, coordinate descent)
assign stable identities but are outside this paper's scope.

\paragraph{Integrity scope.} We enforce parent-integrity (present $\Rightarrow$
parent present). Sibling-order integrity (no duplicate \texttt{order} after
merge) and deeper structural invariants (uniqueness of the root, acyclicity
of $parent$) are natural further sheaf constraints left to future work.

\paragraph{Snapshot vs.\ patch commutation.} Our merge is snapshot-based
(pushout on sections), not patch-based (commutation on operations).
Cherry-pick and rebase are \texttt{apply}-based and less flexible than
Pijul's. A full commutation algebra for tree patches is future work.

\paragraph{Empirical evaluation.} A real-world study on Git repositories
with JSON/YAML, measuring how often the structural-conflict class arises
in practice, would strengthen the practical claims.


% ============================================================
\section{Conclusion}
% ============================================================

We placed the topology on the tree. Merge for structured data, viewed as
a pushout in the sheaf category on the Alexandrov space of ancestry, is
a two-stage computation: a pointwise per-field pushout (the classical
discrete merge) followed by sheafification that enforces referential
integrity. The second stage surfaces a class of obstructions --- a
present child whose parent was concurrently deleted --- that the discrete
per-field merge is provably blind to (Theorem~\ref{thm:strict}). Combined
with identity-location separation, Move + Modify still merges cleanly
because the two edits touch disjoint fields and the moved node's new
parent is present. On a flat tree the sheaf specialises to the discrete
merge; the new content is exactly what the ancestry topology buys.


% ============================================================
\bibliographystyle{plain}
\begin{thebibliography}{9}

\bibitem{semistructured}
S.~Apel, J.~Liebig, B.~Brandl, and C.~Lengauer.
\newblock Semistructured merge: Rethinking merge in revision control systems.
\newblock In \emph{Proc.\ ESEC/FSE}, 2011.

\bibitem{pijul}
P.-{\'E}.~Meunier and contributors.
\newblock Pijul: A sound and fast distributed version control system.
\newblock \url{https://pijul.org}, 2024.

\bibitem{crdt}
M.~Shapiro, N.~Pregui{\c{c}}a, C.~Baquero, and M.~Zawirski.
\newblock Conflict-free replicated data types.
\newblock In \emph{SSS}, 2011.

\bibitem{git}
L.~Torvalds and the Git community.
\newblock Git: A distributed version control system.
\newblock \url{https://git-scm.com}, 2024.

\bibitem{godement}
R.~Godement.
\newblock \emph{Topologie alg\'ebrique et th\'eorie des faisceaux}.
\newblock Hermann, Paris, 1958.

\bibitem{mccord}
M.~C.~McCord.
\newblock Singular homology groups and homotopy types of finite topological
spaces.
\newblock \emph{Topology}, 5(3):249--271, 1966.

\bibitem{barmak}
J.~A.~Barmak.
\newblock \emph{Algebraic Topology of Finite Topological Spaces and
Applications}.
\newblock Lecture Notes in Mathematics 2032, Springer, 2011.

\end{thebibliography}

\end{document}
